\section{Результаты}
\label{sec:results}

\subsection{Общее описание INTENSE}
\label{sec:overview_intense}

Фреймворк INTENSE
предназначен для выявления нейронов, кодирующих определённые поведенческие или средовые переменные, на основе данных крупномасштабного кальциевого имиджинга (\fig{fig:method}). INTENSE работает непосредственно с сырыми сигналами кальциевой флуоресценции, что устраняет необходимость деконволюции спайков. В качестве основной меры используется
взаимная информация (ВИ) --- величина, которая улавливает любую статистическую зависимость, линейную или нелинейную, и обращается в нуль тогда и только тогда, когда переменные независимы
\cite{Shannon1948, Borst1999, Strong1998}. В частности,
в INTENSE применяется оценщик ВИ на основе гауссовой копулы (GCMI)
\cite{gcmi, ma2011mutual}, обеспечивающий эффективное вычисление ВИ в замкнутой форме
(см. Методы: раздел~\ref{sec:mi_computation}), что делает возможным анализ
тысяч пар «нейрон--поведение» в крупных популяциях.

Статистическая значимость оценивается
с помощью пермутационного тестирования на основе циклического сдвига с параметрическим моделированием p-значений
(\fig{fig:method}E; см. Методы: раздел~\ref{sec:statistical_assessment}).
Систематическая оптимизация задержки позволяет учесть кинетику кальциевого индикатора, а также проспективное или ретроспективное кодирование
(раздел~\ref{sec:temporal_alignment}).

\begin{figure}
  \centering
  \includegraphics[width=1\linewidth]{figs/fig1_delegated.png}
  \caption{\textbf{A}: Примеры нормированных следов кальциевой флуоресценции $dF/F$ (синий) с детектированными событиями (серый).
\textbf{B}: Примеры поведенческих/средовых
переменных, используемых для корреляции с нейронной активностью: непрерывная (направление головы, вверху),
дискретная (взаимодействие с объектом, середина), многомерная (координаты животного,
внизу). \textbf{C}: Схема конвейера INTENSE для количественной оценки нейронной селективности. \textbf{D}: Слева: пример клетки, селективной к взаимодействию с объектом:
след флуоресценции (синий) и периоды взаимодействия (зелёный). Справа: распределения нормированных значений $dF/F$ во время (зелёный) и вне (серый) периодов взаимодействия с объектом. \textbf{E}: Распределение
пар «нейрон--поведение» в пространстве «значимость--сила эффекта» с
маргинальными распределениями. Пары, признанные значимыми, показаны зелёным.}

  \label{fig:method}
\end{figure}

Помимо детекции, INTENSE предоставляет инструменты для анализа смешанной селективности
(\fig{fig:results}F,G). Многие нейроны гиппокампа и коры демонстрируют видимую
селективность к нескольким поведенческим переменным одновременно \cite{mix2,mix3}, однако отделение
истинной смешанной селективности от корреляций, индуцированных ковариирующими поведенческими переменными, имеет принципиальное значение
для понимания принципов нейронного кодирования и истинной размерности
представлений. Процедура разделения запутанных селективностей, реализованная в фреймворке, использует условную ВИ и информацию взаимодействия \cite{ghassami2017} для отделения подлинного многопеременного кодирования от
ложных ассоциаций (раздел~\ref{sec:disentangling_selectivity},
раздел~\ref{sec:intense_deciphers}).

Валидация INTENSE была проведена на синтетических наборах данных с известным истинным ответом в широком
диапазоне отношений сигнал/шум (SNR) и условий надёжности ответа
(\fig{fig:synthetics}, раздел~\ref{sec:synthetic_validation}), а также путём сопоставления
с классическими методами детекции клеток места (КМ), основанными на событиях
(\fig{fig:pc}, раздел~\ref{sec:intense_validates}). Затем фреймворк был применён к данным
кальциевого имиджинга гиппокампа свободно перемещающихся мышей, что позволило выявить разнообразные
функциональные типы клеток, включая КМ, клетки направления головы и нейроны с
подлинной смешанной селективностью (\fig{fig:results}A,B,D,E,
раздел~\ref{sec:intense_deciphers}).


\subsection{Валидация INTENSE сопоставлением с устоявшимися методами детекции клеток места}
\label{sec:intense_validates}

\begin{figure}
  \centering
  \includegraphics[width=0.9\linewidth]{figs/fig2_delegated.png}
  \caption{\textbf{A}: Вверху --- карта непрерывного вейвлет-преобразования (CWT) сигнала
флуоресценции. Гребни, проходящие через несколько масштабов, соответствуют детектированным событиям. Внизу --- след флуоресценции (синий) и детектированные события (серый). \textbf{B}: Пример траектории животного, цветом показана нормированная кальциевая активность одной клетки места. Звёздочками отмечены детектированные события. \textbf{C}: Слева --- пересечение
популяций КМ, идентифицированных INTENSE (красный) и классическим методом (синий) при
$p<0.05$. Справа: то же для всех клеток, идентифицированных как КМ любым из методов. \textbf{D}: Пример карт активности, построенных по
сырым кальциевым сигналам (слева) и по детектированным событиям (справа) с высоким показателем SSIM.
\textbf{E}: Метрики сходства между картами активности, построенными по сырым кальциевым
сигналам и по детектированным событиям, как функции уровня доверия (измеренного через
$p$-значение). Слева: среднеквадратичная ошибка (MSE); в центре: пиковое отношение сигнал/шум (PSNR); справа:
индекс структурного сходства (SSIM). Затенённые области соответствуют доверительным интервалам.
Вертикальными линиями показан порог $p=0.001$, как и на \fig{fig:pc}F. \textbf{F}: Относительное пересечение популяций КМ, выявленных
INTENSE и классическим методом, как функция уровня доверия (измеренного через
$p$-значение). Пересечение вычислено для клеток, у которых $p_{INTENSE} \leq p$ и
$p_{classic} \leq p$ одновременно. Нейроны, классифицированные как «не КМ» обоими
методами, исключены. Показаны пересечения как доли соответствующих
популяций (красный для INTENSE, синий для классического метода); коэффициент Жаккара
$J$=(INTENSE $\bigcap$ Classic) / (INTENSE $\bigcup$ Classic), оранжевый; доля
всех идентифицированных КМ (любым методом), оставшихся после порогового отсечения по $p$.}
  \label{fig:pc}
\end{figure}


Для проверки точности INTENSE фреймворк был применён к хорошо изученной задаче
идентификации клеток места (КМ) --- нейронов, проявляющих селективность к определённым местоположениям в
окружающей среде \cite{o1971hippocampus}.

Детекция КМ представляет собой требовательный критерий валидации, поскольку выбор
метода классификации в значительной степени определяет выявляемую популяцию клеток
\cite{Grijseels2021}: перекрытие между различными методами, основанными на событиях, может составлять всего 30--40\%.

Было проведено сравнение INTENSE, применённого непосредственно к сигналам кальциевой флуоресценции и
координатам животного (см. Методы: раздел~\ref{sec:intense_pc_identification}), с
классическим методом на основе событий, строящим бинаризованные карты активности по деконвольвированным
кальциевым событиям \cite{skaggs1993information} (см. Методы:
раздел~\ref{sec:classic_pc_identification}). Для обеспечения прямого и корректного сопоставления результаты обоих методов были приведены к единому формату: (i) бинарная метка КМ и (ii) соответствующий показатель уверенности (p-значение), полученный на основе собственной нулевой модели каждого метода (см.
Методы: раздел~\ref{sec:comparison_pc_populations}).

Несмотря на умеренное перекрытие ($\sim$40\%) при учёте всех клеток, идентифицированных любым из
методов, согласованность резко возрастала при применении порогов уверенности
(\fig{fig:pc}C). Перекрытие монотонно увеличивалось с ужесточением порога, достигая
90\% при $p < 0.001$ --- пороге, при котором сохранялось 95\% общей популяции КМ
(\fig{fig:pc}F). Дальнейшее повышение порога приводило к почти полному
согласию, однако ценой исключения значительной части предполагаемых КМ.

Существенная гетерогенность частот разрядов и стабильности КМ установлена как
экспериментально \cite{Mizuseki2012}, так и теоретически \cite{Lian2022}.
Полученные результаты свидетельствуют о том, что INTENSE и методы, основанные на спайках, различаются главным образом в отношении
«слабых» КМ, но сходятся на одной и той же «ядерной» популяции при повышении
требований к статистической достоверности. Таким образом, предложенный общий
информационно-теоретический подход успешно выявляет свойства пространственного кодирования,
идентифицируемые специализированными алгоритмами.


Для подтверждения того, что все совпадающие КМ, идентифицированные обоими методами, обладают одинаковым паттерном
пространственной селективности, был проведён анализ совпадения карт активности. Существует
широкий спектр методов детекции полей места \cite{Grieves2017}, но они в
значительной степени зависят от типа данных. Для преодоления этой трудности было проведено сравнение
дискретизированных карт активности (нормированных на время пребывания в ячейках; см. Методы:
раздел~\ref{sec:comparison_spatial_selectivity} и \fig{fig:pc}D).

Соответствие между картами активности было количественно оценено с использованием индекса структурного
сходства (SSIM) и дополнительных метрик (см. Методы:
раздел~\ref{sec:comparison_spatial_selectivity}).

Все метрики пространственного соответствия улучшались с повышением порогов уверенности
(\fig{fig:pc}E), что подтверждает: КМ, идентифицированные с высокой уверенностью обоими
методами, демонстрируют весьма сходные паттерны пространственной активности. Данная многокритериальная
валидация показывает, что INTENSE не только выявляет совпадающие популяции клеток,
но и точно фиксирует их функциональные свойства.


\subsection{INTENSE выявляет разнообразную нейронную селективность}
\label{sec:intense_deciphers}

\begin{figure}
  \centering
  \includegraphics[width=1\linewidth]{figs/fig3_delegated.png}
  \caption{\textbf{A}: Пример «нейрона взаимодействия с объектом» и распределения $dF/F$ во время взаимодействия (зелёный) и в
контрольные периоды (серый). \textbf{B}:
Пример «нейрона вставания на задние лапы» и распределения $dF/F$ во время вставания (оранжевый) и в контрольные периоды (серый). \textbf{C}: Слева --- распределение нейронов, селективных к 1--3 переменным. Справа --- распределение оптимальных задержек: отрицательные ($-$), околонулевые и положительные
($+$). Статистические сравнения подробно описаны в
Методах. \textbf{D}: Карта настройки на направление головы мультиселективного нейрона.
\textbf{E}: Двумерная карта пространственной активности того же нейрона. \textbf{F}: Тепловая карта ВИ между
статистически значимо коррелирующими поведенческими признаками (анализ «поведение--поведение»). Сила связи отражена интенсивностью цвета.
\textbf{G}: Схема взаимозависимостей между информационными мерами для
нейронов с запутанной селективностью.}

  \label{fig:results}
\end{figure}


После валидации INTENSE сопоставлением с устоявшимися методами фреймворк был применён для обнаружения
новых ассоциаций между нейронами и поведением в гиппокампальных записях. Были проанализированы
ассоциации между нейронной активностью и широким набором поведенческих признаков,
включающим пространственное положение, направление головы, параметры движения и дискретные
поведенческие состояния (см. Методы: раздел~\ref{sec:behavioral_analysis}). После
идентификации значимых пар «нейрон--поведение» была проведена
процедура разделения запутанных селективностей для разграничения истинной смешанной селективности и ложных
ассоциаций, возникающих вследствие поведенческих корреляций (\fig{fig:results}A--E; см.
Методы: раздел~\ref{sec:disentangling_selectivity}).

Процедура разделения использует условную взаимную информацию (УВИ) для
разграничения первичных и вторичных селективностей, а информацию взаимодействия (ИВ) --- для
классификации пар признаков как избыточных или синергических (см. Методы:
раздел~\ref{sec:disentangling_selectivity}).

Результаты показали наличие множества нейронов, селективных к различным
конкретным поведенческим признакам (таким как взаимодействие с объектом, вставание на задние лапы и др.,
как показано на \fig{fig:results}A,B).
% [AUTHOR: The following paragraph on probabilistic activation patterns needs review.
% Consider whether this interpretive discussion belongs in Results or Discussion.]
Эти поведенчески селективные нейроны демонстрировали
стохастический характер активации --- они не разряжались детерминированно при
каждом эпизоде предпочтительного поведения и не оставались полностью неактивными в остальное время.
Это может быть объяснено вероятностной природой нейронного кодирования и необходимостью
более детальной идентификации поведенческих актов с учётом поведенческого контекста.
Акты, которые могут казаться идентичными с точки зрения исследователя на основании
кинематических характеристик реперных точек, могут нести различный внутренний
смысл для животного.

Наряду с кодированием дискретных поведенческих актов были обнаружены многочисленные клетки, селективные к
непрерывным переменным: скорости движения, направлению головы и пространственному положению
(\fig{fig:results}D,E). Клетки, селективные к состоянию локомоции, направлению головы или скорости, всё чаще выявляются в
гиппокампе \cite{kropff2015speed,Gois2018,Acharya2016}.

\paragraph{Временное согласование нейронной активности и поведения.}
\label{sec:temporal_results}
Для учёта временных рассогласований между нейронной активностью и сигналом кальциевой
флуоресценции в INTENSE проводится оптимизация задержки между сигналами (см. Методы:
раздел~\ref{sec:temporal_alignment}).
Оптимальные задержки были разделены на три категории (отрицательные, околонулевые
и положительные) для статистического анализа. Большинство значимых ассоциаций
наблюдалось при положительных задержках (44,5\%), что согласуется с ожидаемым временным запаздыванием
кальциевых индикаторов; 41,0\% --- при околонулевых задержках и лишь 14,5\% --- при отрицательных
(\fig{fig:results}C). Поскольку основной целью являлась характеристика ассоциаций, согласующихся с запаздыванием кальциевого индикатора, распределения задержек приводятся для всех значимых пар, а последующие анализы селективности, если не указано иное, ограничены парами с неотрицательными оптимальными задержками.

\paragraph{Количественная оценка и разделение смешанной селективности.}
Из 30{,}105 пар «нейрон--сессия», проанализированных по 16 мышам и 4 сессиям, 6{,}217
(20,7\%) показали значимую селективность хотя бы к одной поведенческой переменной.
Среди них 1{,}377 (22,1\% селективных нейронов) отвечали на две и более
переменные до проведения процедуры разделения (\fig{fig:results}C).

Попарный анализ значимости поведенческих переменных выявил тесно связанные
кластеры (\fig{fig:results}F). Переменные, связанные с локомоцией (скорость, покой,
ходьба, состояние локомоции), показали значимую ко-селективность в 82--100\%
комбинаций «мышь--сессия», что указывает на сильные поведенческие корреляции, способные
завышать оценку смешанной селективности. Аналогичным образом переменные пространственных зон (углы,
стены, центр) сформировали отдельный кластер (77--91\% ко-значимости).

Процедура разделения на основе УВИ позволила разрешить существенную долю этих
многопеременных ассоциаций. После контроля поведенческих корреляций
число мультиселективных нейронов уменьшилось с 1{,}377 до 949 (15,7\% от
селективных); примерно треть многопеременных ассоциаций была объяснена
поведенческой избыточностью. Наиболее выраженный эффект наблюдался для селективности к скорости: из 201
нейрона, первоначально классифицированного как селективный к скорости, лишь 7 сохранили значимость после
кондиционирования на состояние локомоции, что согласуется со 100\%-ной ко-селективностью скорости и локомоции во всех сессиях.

Информация взаимодействия обеспечила дополнительную характеристику ассоциаций поведенческих переменных.
Для пары «скорость--локомоция» ИВ была устойчиво отрицательной (избыточность), что подтверждает
перекрывающуюся информацию этих переменных о нейронной активности.
Напротив, такие пары, как направление тела и направление головы, демонстрировали близкую к нулю или
положительную ИВ в большинстве сессий, что указывает на то, что эти тесно связанные переменные
несут в значительной степени совпадающую информацию о нейронной активности.

Даже после проведения процедуры разделения значительное число нейронов сохраняло значимые
ассоциации с несколькими переменными. Например, нейрон, показанный на
\fig{fig:results}D,E, демонстрировал подлинную смешанную селективность к скорости, направлению головы и
пространственному положению одновременно. Хотя такие мультиселективные нейроны составляют меньшинство
в популяции (\fig{fig:results}C), они могут играть ключевую роль в интеграции
различных функциональных систем \cite{mix3}.

\paragraph{Взаимная информация выявляет ассоциации, невидимые для линейной корреляции.}
\label{sec:mi_vs_corr}
При замене ВИ на корреляцию Пирсона в идентичном во всём остальном конвейере INTENSE было выявлено в 2,2 раза меньше селективных нейронов и в 2,9 раза меньше значимых пар «нейрон--признак». Наибольшее расхождение наблюдалось для дискретных поведенческих переменных (5--11$\times$) и оставалось существенным даже для непрерывных переменных, таких как скорость (3,2$\times$). Эти результаты подтверждают, что безмодельная информационно-теоретическая метрика существенно расширяет репертуар детектируемых нейронных селективностей.

\subsection{Валидация на синтетических данных демонстрирует робастность INTENSE}
\label{sec:synthetic_validation}

\begin{figure}
  \centering
  \includegraphics[width=1\linewidth]{figs/fig4_delegated.png}
  \caption{\textbf{A}: Этапы генерации псевдокальциевого сигнала, связанного с
непрерывной переменной. 1 --- непрерывный поведенческий временной ряд с выделенным диапазоном значений;
2 --- рандомизированный бинарный временной ряд (с $p_{skip}$); 3 --- неоднородный
пуассоновский процесс с двумя различными интенсивностями; 4 --- итоговый псевдокальциевый сигнал после
свёртки с ядром кальциевого индикатора. Для дискретных переменных процедура
начинается с этапа~2.
\textbf{B}: Тепловые карты $F_1$-меры для теста с непрерывной переменной во всём
диапазоне параметров. Показаны методы: на основе корреляции, на основе ВИ, на основе корреляции с
перемешиваниями, INTENSE. \textbf{C}: То же для дискретных переменных (сгенерированных начиная
с этапа~2 из A). Показаны методы: на основе среднего, на основе ВИ,
на основе среднего с перемешиваниями, INTENSE. \textbf{D}: Карты «точность--полнота» для
теста с непрерывной переменной. Показаны методы: на основе корреляции, на основе ВИ,
на основе корреляции с перемешиваниями, INTENSE. Слева --- координаты «точность--полнота»
всех методов для различных значений SNR (при $p_{skip} = 0$). Справа ---
координаты «точность--полнота» всех методов для различных значений $p_{skip}$ (при
$SNR = 64$). Вертикальными линиями показана точность случайного угадывания ($PR_{random}
=0.05$). \textbf{E}: Карты «точность--полнота» для теста с дискретной переменной. Показаны
методы: на основе среднего, на основе ВИ, на основе среднего с перемешиваниями, INTENSE. Слева ---
координаты «точность--полнота» всех методов для различных значений SNR (при
$p_{skip} = 0$). Справа --- координаты «точность--полнота» всех методов для различных
значений $p_{skip}$ (при $SNR = 64$). Вертикальными линиями показана точность случайного
угадывания ($PR_{random} =0.05$).}
  \label{fig:synthetics}

\end{figure}


Для строгой оценки детекционных возможностей INTENSE были проведены систематические
испытания на синтетических наборах данных с известным истинным ответом. Генерировались эксперименты,
содержащие 20 поведенческих временных рядов и 500 псевдокальциевых сигналов, причём каждый
кальциевый сигнал был связан ровно с одной поведенческой переменной (распространённость истинных
ассоциаций составляла 5\%).

\subsubsection{Синтетическая модель и роль контролей временной структуры}
\label{sec:bio_synthetic}

Были сгенерированы реалистичные синтетические эксперименты, в которых кальциевые сигналы были
связаны с поведенческими переменными через биологически правдоподобный процесс
(\fig{fig:synthetics}A; см. Методы: раздел~\ref{sec:synthetic_data_generation}).
Варьировались отношение сигнал/шум
(SNR $\in \{2, 4, 8, 16, 32, 64\}$) и надёжность ответа
($p_{\text{skip}} \in [0.0, 0.8]$).

Сравнение методов выявило фундаментальную проблему: временная автокорреляция
порождает подавляющее число ложноположительных результатов в отсутствие надлежащего статистического контроля
(\fig{fig:synthetics}B--E). Методы, использующие стандартные тесты на основе корреляции или ВИ
без контроля перемешиванием, достигали точности лишь 8,7\% и 6,4\%
соответственно для непрерывных переменных, что едва превышает 5\%-ный уровень случайного угадывания.
Это демонстрирует, почему пермутационное тестирование с перемешиванием необходимо --- медленная
динамика как кальциевых сигналов, так и поведенческих переменных порождает ложные
корреляции, которые выглядят высокозначимыми в рамках наивных статистических тестов.

\subsubsection{INTENSE превосходит другие подходы с перемешиванием}
\label{sec:intense_outperforms}

Среди методов, использующих контроль перемешиванием, INTENSE продемонстрировал явные преимущества перед
подходами на основе корреляции. Хотя метод на основе корреляции с контролем перемешиванием достиг высокой
точности (100,0\% при SNR = 64), его полнота была крайне низкой
(максимум 4,5\%) ---
обнаруживалось менее 5\% истинных ассоциаций даже в оптимальных
условиях. Этот паттерн указывает на принципиальное ограничение: методы на основе корреляции,
даже при надлежащем статистическом контроле, ограничены линейными
ассоциациями и не улавливают сложные нелинейные зависимости, часто присутствующие в нейронных
данных.

Напротив, INTENSE показал устойчивую работу во всём спектре
интенсивностей сигнала. При идеальной надёжности ответа ($p_{\text{skip}} = 0$)
точность возрастала с 44,2\% при SNR = 2 до 85,5\% при SNR = 64, тогда как полнота
улучшалась с 6,8\% до 41,7\% для непрерывных переменных.
Для дискретных поведенческих переменных были получены существенно лучшие результаты: точность возросла с 73,5\% до
94,6\%, а полнота увеличилась с 24,3\% до 100,0\% в том же
диапазоне.

\subsubsection{Двухэтапное тестирование обеспечивает баланс между эффективностью и мощностью}
\label{sec:two_stage_testing}

Двухэтапная процедура тестирования (см. Методы: раздел~\ref{sec:statistical_assessment})
показала высокую эффективность в управлении
вычислительной нагрузкой крупномасштабного анализа. На первом этапе отбиралось 3--6\%
всех возможных пар в качестве кандидатов.
Точность возрастала с 62,0\% до 85,5\% для непрерывных переменных и с 89,4\% до 94,6\%
для дискретных между этапами, тогда как полнота снижалась лишь незначительно, что свидетельствует об
эффективном контроле ложных обнаружений.

\subsubsection{Работа в условиях биологической вариативности}
\label{sec:performance_variability}

Реальные нейроны демонстрируют непостоянную надёжность ответа --- они могут пропускать ответы вследствие
адаптации, изменений внутреннего состояния или стохастических факторов. INTENSE продемонстрировал
плавную деградацию характеристик при снижении надёжности ответа. При SNR = 64
точность для непрерывных переменных снижалась с 85,5\% при $p_{\text{skip}} = 0$ до
74,1\% при $p_{\text{skip}} = 0,5$. Дискретные переменные показали большую устойчивость:
точность сохранялась на уровне 94,8\% при полноте 99,7\% даже при $p_{\text{skip}} = 0,5$,
что свидетельствует о робастности к умеренной ненадёжности ответа, типичной для реальных нейронных данных.

Любопытный эффект проявился при очень высоких отношениях сигнал/шум: метрики на основе сырой ВИ
и средних значений демонстрировали ухудшение характеристик вследствие эффектов насыщения кальциевого
индикатора (\fig{fig:synthetics}B,C). Когда нейроны разряжаются часто
и интенсивно, медленная динамика затухания приводит к тому, что флуоресценция остаётся повышенной между событиями,
сливая отдельные эпизоды в протяжённые плато. Для дискретных переменных это вызвало
снижение точности сырой ВИ с 92,9\% при SNR = 4 до 70,0\% при SNR = 64, а
детекция на основе среднего с контролем перемешиванием деградировала от $F_1$ = 51,9\% до
10,1\% в том же диапазоне. INTENSE оказался устойчив к этим эффектам: $F_1$-мера монотонно возрастала
с 95,2\% до 97,2\%. Данная робастность обусловлена тем, что циклические
сдвиги строят нулевое распределение из тех же насыщенных сигналов, поэтому насыщение
повышает как истинные, так и нулевые значения ВИ в равной мере, сохраняя валидность
теста значимости.

Результаты синтетической валидации демонстрируют, что INTENSE успешно решает ключевые
задачи детекции ассоциаций «нейрон--поведение»: контроль ложных обнаружений,
обусловленных временной структурой данных, выявление сложных ассоциаций за пределами линейной корреляции
и сохранение чувствительности в условиях биологической вариативности. Робастная
работа фреймворка в разнообразных условиях обосновывает уверенность в его применимости к экспериментальным данным, для которых истинный ответ неизвестен.
